% ============================================================
%  Chapter: Attitude, Deontology and Utilitarianism
%  Class Notes L11-12
%  Topics: Attitude \textbullet\ Deontology (Immanuel Kant)
%          \textbullet\ Utilitarianism (Bentham \& Mill)
% ============================================================

\chapter{Attitude, Deontology and Utilitarianism}

\noindent{\sffamily\color{Grey}\fontsize{9.5}{12}\selectfont
Class Notes L11-12\\
Attitude \textbullet\ Deontology (Immanuel Kant) \textbullet\ Utilitarianism (Bentham \& Mill)}
\vspace{14pt}

\section{Attitude (continued)}

\subsection{Cognitive Dissonance}

\begin{KeyConcept}
\begin{itemize}
\item Cognitive dissonance is the mental discomfort that results from holding two conflicting beliefs, values or attitudes.
\item It arises when what we think and what we do are not in harmony, or when our attitude and our behaviour are not in harmony (incongruence between attitude and behaviour).
\item This incongruence produces mental discomfort, which in turn creates problems for the individual.
\end{itemize}
\end{KeyConcept}

Cognition simply means a belief. When our behaviour contradicts a belief we hold, dissonance is generated.

\begin{ExampleBox}
A person has a strong positive attitude towards smoking, but also holds the belief that smoking causes cancer. The clash between the two produces mental discomfort --- this is cognitive dissonance.
\end{ExampleBox}

\paragraph{Ways to reduce cognitive dissonance}
\begin{itemize}
\item Change the behaviour: for example, quit smoking so that the behaviour aligns with the belief.
\item Change / add to the belief pattern: justify the behaviour by adding a new belief --- e.g. ``I am allowed to have some cheat meals every once in a while.''
\item Justify or deny the conflicting cognition so that the conflict is neutralised.
\end{itemize}

In every case, dissonance is reduced by changing behaviour, or by changing the conflicting cognition that is in conflict with the behaviour.

\begin{QuickRevision}
\begin{itemize}
\item Cognitive dissonance = mental discomfort from two conflicting beliefs / values / attitudes (attitude--behaviour incongruence).
\item Reduced by (a) changing behaviour, or (b) changing / justifying / denying the conflicting cognition.
\end{itemize}
\end{QuickRevision}

\subsection{Political Attitude}

\begin{KeyConcept}
\begin{itemize}
\item Attitude is a predisposition to act in a certain way. Predisposition means our tendency to act towards an attitudinal object.
\item In political attitude the attitudinal object is fixed as politics --- a political party, a leader, an ideology or a policy --- and the question is whether our tendency is to support it or to oppose it.
\item It is therefore an expression of favour or disfavour towards particular ideologies, political parties, political culture and politicians --- essentially everything concerning politics.
\end{itemize}
\end{KeyConcept}

Evaluation here means favour or disfavour (e.g. a positive or negative attitude towards a leader). Attitude is a relational term; in political attitude that relation is established with politics.

\paragraph{Factors that shape political attitude}

\paragraph{Historical factors}
\begin{itemize}
\item Historical experiences such as the freedom struggle taught us the importance of sovereignty --- that all of us must safeguard it.
\item Having lived through colonial subjugation under the British East India Company, if asked today to choose between the comfort of aid from a developed nation and sovereignty, most of us would choose sovereignty. History thus shapes our political attitude towards sovereignty.
\end{itemize}

\paragraph{Social factors --- caste}
\begin{itemize}
\item Caste-based societies and caste-based political parties mean our attitude is often shaped by our own caste identity. As it is said: ``In India, you don't cast your vote, you vote your caste.''
\item Many political parties have emerged from caste-based politics; if a leader belongs to your caste, you are most likely to favour that party.
\end{itemize}

\paragraph{Gender}
\begin{itemize}
\item Women may vote for parties that favour women's agendas or issues.
\end{itemize}

\begin{GovtPolicy}
Example --- Nitish Kumar (then Chief Minister of Bihar) promised liquor prohibition (sharab-bandi), addressing the problems of women subjected to domestic violence by alcoholic husbands. This women-centric policy helped form a positive political attitude towards his party.
\end{GovtPolicy}

\paragraph{Family}
\begin{itemize}
\item People say ``we are Congress people'' or ``we are BJP people.'' Generation after generation a family votes for one party / ideology, so an individual's political attitude tends to remain the same.
\end{itemize}

\paragraph{Education}
\begin{itemize}
\item As education increases, a person understands rights better, becomes more liberal, grows more conscious of their rights, and demands them from government --- protesting or posting on social media. All of this is an expression of political attitude.
\end{itemize}

\paragraph{Age}
\begin{itemize}
\item Young people tend to believe in total revolution and rapid change; older, more mature people favour incremental reforms, understanding that real change on the ground happens slowly. Age therefore shapes the maturity of one's political attitude.
\end{itemize}

\paragraph{Economic status}
\begin{itemize}
\item A poor person who benefits from government subsidies favours parties aligned with subsidies and welfare programmes; a rich person may favour parties offering tax breaks and reductions and dislike subsidies.
\end{itemize}

\paragraph{Media}
\begin{itemize}
\item Social media and media narratives, through campaigns and advertisements, are used to build a narrative or a positive attitude towards a particular political party.
\end{itemize}

\paragraph{Circumstances}

\begin{ExampleBox}
\begin{itemize}
\item The rise of Hitler in Germany was significantly shaped by circumstances. His party won more votes during economic uncertainty --- after the 1929 Great Depression, Germany faced very high inflation, unemployment and severe economic challenges, and in that context people looked to Hitler to solve their problems.
\item In dire circumstances people seek massive change, and charismatic leaders who promise it are able to influence them.
\end{itemize}
\end{ExampleBox}

Additional factors (from the PDF): educational institutions, leaders in society, and political groups themselves also play a large role in shaping political attitude. This can appear as an independent question in UPSC.

\begin{QuickRevision}
\begin{itemize}
\item Political attitude = favour / disfavour towards parties, ideologies, leaders and policies.
\item Shaping factors: historical, caste, gender, family, education, age, economic status, media, circumstances (plus educational institutions, societal leaders, political groups).
\end{itemize}
\end{QuickRevision}

\subsection{Social Influence}

\begin{KeyConcept}
\begin{itemize}
\item Social influence is the change in our behaviour or attitude caused by an external / actual factor around us (for instance, ``this boy fell into bad company'' --- his attitude and behaviour changed wrongly because of his company).
\item The external pressure may be real or imaginary.
\item The influence may be intentional or unintentional.
\end{itemize}
\end{KeyConcept}

\begin{ExampleBox}
\begin{itemize}
\item A celebrity such as Ranveer Singh dressing in a certain way may not himself know how many people he influences --- the influence is unintentional. Social-media influencers are similar.
\item In the Samay Raina controversy, the concern was that such content may adversely impact children; mature people can process it more maturely, but on children it may have a negative influence.
\item A child enters college, sees others smoking, and starts smoking to look `cool' --- no one told him to, he was simply influenced by the people around him (a real or imaginary pressure).
\end{itemize}
\end{ExampleBox}

Bollywood stars, cricketers, celebrities, social-media influencers and the people immediately around us all influence our behaviour and attitude.

\paragraph{How social influence operates}

\paragraph{Informational influence}
\begin{itemize}
\item Social influence is more operative when it comes from a source of credibility.
\end{itemize}

\begin{ExampleBox}
\begin{itemize}
\item During COVID, a reputed doctor held a press conference urging people to stop getting repeated CT scans --- too many scans were taking a toll on medical infrastructure, so those who genuinely needed them could not get them. When a credible, professional person speaks, the influence is stronger.
\item IPCC reports similarly provide informational influence about the consequences of climate change and global warming (such as sea-level rise).
\end{itemize}
\end{ExampleBox}

\paragraph{Situational influence}
\begin{itemize}
\item The situation a person is in also influences their attitude and behaviour.
\end{itemize}

\paragraph{Normative influence}
\begin{itemize}
\item Norms influence us --- society expects certain behaviour, e.g. that we respect our elders. When you talk to your parents knowing that everyone is watching, the (normative) influence increases; when onlookers do not know the parents, that normative influence does not operate as strongly.
\end{itemize}

\paragraph{Factors determining the degree of social influence}
\begin{itemize}
\item Group size: a larger group exerts more influence.
\item Cohesiveness: the more cohesive the group, the greater its influence.
\item Importance: if a matter is highly important for one's individual decision (e.g. one's own marriage), the person will not be easily influenced even within a large group.
\item Power distance / power equation: if the power distance is very low we are not influenced easily; a very powerful person (e.g. a celebrity, the PM, a DM) has high power distance and can influence us more easily.
\item Individual (moral) character: a person of strong moral character is less easily influenced.
\end{itemize}

\begin{CriticalPoint}
Distinguish undue influence from coercion. ``Everyone is doing it, so I did it too'' is undue influence --- a negative example of social influence.
\end{CriticalPoint}

\begin{QuickRevision}
\begin{itemize}
\item Social influence = external factor (real or imaginary; intentional or unintentional) changing attitude / behaviour.
\item Operates through: informational, situational and normative influence.
\item Degree depends on: group size, cohesiveness, importance, power distance, individual character.
\end{itemize}
\end{QuickRevision}

\subsection{Persuasion}

\begin{KeyConcept}
\begin{itemize}
\item Persuasion involves an intentional communicative act that excludes force and achieves private acceptance.
\item It is deliberate and intentional --- the purpose is to intentionally change someone's attitude or behaviour.
\end{itemize}
\end{KeyConcept}

\begin{CriticalPoint}
Persuasion is a specialised case of social influence: every persuasion is a case of influence, but not every influence is a case of persuasion. In social influence the change need not be intentional (celebrities may or may not intend to influence you); in persuasion we truly and deliberately intend to influence.
\end{CriticalPoint}

\paragraph{Techniques of persuasion}

\paragraph{(a) Social proofing} Showing what other people are doing in order to influence a person --- telling them ``you are not alone, many others have done it,'' so they too should change their attitude or behaviour.

\begin{GovtPolicy}
\begin{itemize}
\item PM CARES Fund: ``people from all walks of life have contributed \ldots\ they have given their heart and money to sharpen the fight against COVID-19''; celebrity names (e.g. Badshah, Guru Randhawa) were added.
\item `Give it up' subsidy campaign: telling people that one crore citizens have already given up their subsidy, so they are not alone in doing so.
\item Selfie-with-Daughter campaign: people post pictures showing pride in their daughters --- signalling that a daughter is a matter of pride, not shame.
\item `My Waste, My Responsibility' campaign.
\end{itemize}
\end{GovtPolicy}

\paragraph{(b) Default rule} Humans are by nature status-quoist --- we generally do not change things. Making something the default therefore changes attitude and behaviour. This is why every app asks you to `make it default.'

\begin{DataFact}
Austria treats every individual as an organ donor at the time of death by default; those who do not wish to donate can follow a simple opt-out process. Result: more than 90\% of people are registered to donate their organs.
\end{DataFact}

\begin{GovtPolicy}
\begin{itemize}
\item Government uses the auto-debit facility to enrol people in insurance by default (automatically enrolled, can withdraw anytime).
\item TDS (Tax Deducted at Source): the employer deducts a set percentage of salary and deposits it with the tax department by default; if the tax liability does not arise, the person files an ITR and claims the amount back. The default is TDS.
\item Aadhaar was made the default in many services, which made it easy to change people's behaviour --- they began using Aadhaar everywhere.
\end{itemize}
\end{GovtPolicy}

\paragraph{(c) Scarcity complex} Showing that something is limited (``only a few seats left'') creates a FOMO effect and nudges quick action. Used by organisations and government (limited free seats, fees thereafter); e-commerce sites (Myntra, Flipkart) show ``only one / two items left.''

\paragraph{(d) Reciprocity} Creating a moral debt --- give something first, build trust, and the other person then feels obligated to reciprocate (e.g. by buying your product).

\begin{ExampleBox}
\begin{itemize}
\item Jio gave free internet initially; people got used to it and later paid.
\item On roadways buses, a person applies balm on your forehead unasked, creating a moral debt so you feel you should give something in return.
\item Coaching institutes offer free open mock tests --- doing something for you first to create obligation.
\end{itemize}
\end{ExampleBox}

\paragraph{(e) Attention \& clear messaging} A very clear, targeted message changes attitude.

\begin{GovtPolicy}
\begin{itemize}
\item `Do boond zindagi ki' (two drops of life) --- the polio vaccine is equated with life itself.
\item `Give it up, feel the joy of giving' --- giving is equated with joy.
\item Namami Gange: the Ganga is addressed as `Ma' (mother) and `Namami' means we bow to her --- so how can we pollute her? The name itself carries the message (as do names like Mudra Yojana).
\end{itemize}
\end{GovtPolicy}

\paragraph{(f) Messenger liking} Using a liked celebrity or brand ambassador so that people respond to the messenger more than the message.

\begin{GovtPolicy}
\begin{itemize}
\item The Election Commission of India appoints cricketers as brand ambassadors during elections; MS Dhoni appeals to Jharkhand, Ganguly to Bengal, Sachin Tendulkar broadly. `Pehle matdaan, phir jalpaan' (first vote, then refreshments) links voting to responsibility.
\item People buy Patanjali products because they like Baba Ramdev. Kailash Satyarthi is used for child-welfare messaging because he evokes respect.
\item Swachh Bharat Mission engaged many celebrities --- a classic case of persuasion and attitudinal change.
\end{itemize}
\end{GovtPolicy}

\paragraph{(g) Public commitment} Making people publicly commit to something (e.g. at an anti-addiction rally). A public commitment creates guilt if not followed, because people wish to protect their self-image, credibility and trustworthiness.

\begin{ExampleBox}
`Mirror, mirror on the wall' --- mirrors installed at spots where people urinate in public in Bengaluru, to instil public shame and deter the behaviour.
\end{ExampleBox}

\begin{CriticalPoint}
Persuasion can be misused and slip into manipulation (e.g. misleading coaching advertisements). Avoiding manipulation is an important ethical concern.
\end{CriticalPoint}

\begin{QuickRevision}
\begin{itemize}
\item Persuasion = deliberate, intentional attitude / behaviour change without force (a special case of social influence).
\item Techniques: social proofing, default rule, scarcity complex, reciprocity, attention \& clear messaging, messenger liking, public commitment.
\end{itemize}
\end{QuickRevision}

\subsection{Nudging}

Behaviour change lies on a spectrum with two extremes:

\begin{itemize}
\item Laissez-faire --- the choice is left entirely to the individual (e.g. the `Give it up' campaign leaves the decision to each person).
\item Mandate --- the behaviour is compelled by law.
\end{itemize}

Between these two extremes lie nudging and incentives.

\begin{ExampleBox}
Incentive example: a high tax on tobacco raises its cost and thereby reduces tobacco consumption --- behaviour changed through (dis)incentive.
\end{ExampleBox}

\begin{KeyConcept}
Nudging is a `choice architecture' that alters people's behaviour in a predictable way, without changing the choices available and without significantly changing economic incentives. It is a soft push that still respects individual autonomy.
\end{KeyConcept}

\paragraph{Examples of nudging}
\begin{itemize}
\item Metro staircase: the number of kilocalories burnt is printed on the stairs, placed beside an escalator. Both options remain available; the arrangement nudges you to think ``if I take the stairs I'll burn calories.'' No option is removed and there is no mandate.
\item Supermarket shelves: products the store wants you to buy are placed at eye level, while others require you to bend down. All choices remain available; the arrangement nudges the choice.
\item Selfie-with-Daughter, sugar boards, and food-menus with bigger pictures are all nudges; government promotion of EVs is a soft push (nudge) too.
\end{itemize}

\begin{PYQLink}
Nudging is a likely UPSC question --- expect a question on nudging this year or next.
\end{PYQLink}

\begin{CriticalPoint}
Nudging is part of persuasion: like persuasion, it intentionally changes people's attitude --- but it does so through choice architecture rather than force or mandate.
\end{CriticalPoint}

\begin{QuickRevision}
\begin{itemize}
\item Behaviour-change spectrum: laissez-faire $\rightarrow$ nudging $\rightarrow$ incentive $\rightarrow$ mandate.
\item Nudging = choice architecture that predictably alters behaviour without removing choices or significantly changing economic incentives (a soft push).
\end{itemize}
\end{QuickRevision}

\section{Normative Ethics \& Deontology (Immanuel Kant)}

\subsection{Normative Ethics --- The Ethical Standards}

\begin{KeyConcept}
\begin{itemize}
\item Normative ethics deals with prescriptive ethical standards --- what we should consider when deciding the right and wrong of a human action (`norm' = prescription).
\item Different schools prescribe different standards for judging whether an action is ethical.
\end{itemize}
\end{KeyConcept}

\begin{center}
\begin{tabularx}{\textwidth}{>{\raggedright\arraybackslash}p{2.6cm} Y >{\raggedright\arraybackslash}p{2.6cm}}
\rowcolor{AccentDark}
\thead{Approach} & \thead{An action is right when \ldots} & \thead{Focus is on} \\
\toprule
Deontology & you do your duty / the means are right & Duty \& means \\
Teleology & it leads to good consequences / good ends & Ends / consequences \\
Virtue ethics & it makes you a virtuous person & Character of the person \\
Justice & it leads to fairness and justice & Justice / fairness \\
Rights-based & it upholds someone's rights (e.g. a consumer's rights) & Rights \\
\bottomrule
\end{tabularx}
\end{center}

\subsection{Case Study --- SP Sunil (Applying the Standards)}

\begin{ExampleBox}
Sunil is a young, dynamic SP officer posted in a district infamous for illegal sand mining. He begins taking action against the sand mafia; the mafia, rattled, threatens him and his family. One option available to Sunil is to seek a transfer. We evaluate this course of action through the different ethical standards.
\end{ExampleBox}

\begin{center}
\begin{tabularx}{\textwidth}{>{\raggedright\arraybackslash}p{2.4cm} >{\raggedright\arraybackslash}p{2.0cm} Y}
\rowcolor{AccentDark}
\thead{Standard} & \thead{Verdict} & \thead{Reasoning} \\
\toprule
Deontology & Unethical & Seeking a transfer means he is not doing his duty --- he is running away from it. \\
Teleology & Unethical & Bad consequences dominate: more corruption, more illegal sand mining, a bad precedent, poor leadership, and demoralised subordinates (if a senior SP could do nothing). \\
Virtue ethics & Unethical & Seeking a transfer makes him behave like a coward; it does not build good character. \\
\bottomrule
\end{tabularx}
\end{center}

\begin{CriticalPoint}
The same action (seeking transfer) is judged wrong from three different bases --- deontology, teleology and virtue ethics. Examining an action through multiple standards is how you generate dimension in an ethics answer.
\end{CriticalPoint}

\subsection{Kant's Deontology --- Goodwill}

\begin{KeyConcept}
Immanuel Kant: the only thing that is unconditionally / absolutely good is a goodwill. Everything else is a conditional good --- good only depending on how it is used.
\end{KeyConcept}

\paragraph{Everything else is only a conditional good}
\begin{itemize}
\item Courage: a thief, a terrorist and a suicide bomber also need courage --- so courage is a conditional good, depending on how it is used.
\item Love: in the name of love people do terrible things (Dhritarashtra and Gandhari for their son in the Mahabharata; parents willing to pay for leaked exam papers out of love). Love too is a conditional good.
\item Physical strength: it can be misused to beat, threaten or intimidate --- so it is a conditional good. The same applies to beauty and power.
\end{itemize}

\subsection{The Two Features of a Goodwill}

Kant holds that every human being is a rational human being (all of us can think logically). A goodwill has two features.

\paragraph{(i) It must be a Free Will}
\begin{itemize}
\item Free of threat: if a gun is pointed at a person forcing them to steal, they act under threat, not free will.
\item Free of uncontrollable natural urges: a person who has not eaten for fifteen days and steals bread is not acting under free will --- hunger (like the urge to use the washroom) is a natural urge that cannot be controlled.
\item Free of bias / prejudice: if we carry a bias, we are again not operating under free will.
\end{itemize}

\paragraph{Illustration --- the Veil of Ignorance}

\begin{CriticalPoint}
\begin{itemize}
\item The idea of freedom from bias is explained using the Veil of Ignorance. NOTE: this term was given by John Rawls, not by Immanuel Kant --- it is used here in simplified form only to explain Kant's point.
\item `Veil' means a curtain; `ignorance' means we are ignorant about something. Rawls says that whenever we make a decision we should place a veil of ignorance before our eyes --- remaining ignorant of our own identity and interests --- so that we can think rationally and freely.
\end{itemize}
\end{CriticalPoint}

\begin{ExampleBox}
\begin{itemize}
\item Untouchability: eight people in a caste-based society with no government must decide the rules. Under the veil of ignorance they do not know their own caste. Deciding whether untouchability should exist, everyone rationally concludes it should not --- because ``I might turn out to be from the untouchable caste.'' When prejudiced about our own caste, we cannot think freely.
\item Women's education: deciding whether women should be allowed to study, and not knowing their own sex under the veil, everyone concludes women should be allowed to study --- ``because I don't know whether I am male or female.''
\end{itemize}
\end{ExampleBox}

\paragraph{(ii) Follow the Moral Law for the Respect of the Moral Law}
\begin{itemize}
\item A goodwill follows the moral law out of respect for the moral law --- not because of consequences. Consequences mean either reward or punishment.
\end{itemize}

\begin{ExampleBox}
At 12 a.m.\ you are driving and reach a red light --- no CCTV, no policeman, no other car. If you still stop at the red light, knowing there would be no consequence for crossing it, you are following the traffic law for the respect of the traffic law, not to avoid punishment. This is following a moral law for the respect of the moral law.
\end{ExampleBox}

\begin{PYQLink}
\begin{itemize}
\item PYQ (last year): ``In law you are considered guilty when you violate the rights of others; in ethics you are guilty when you even think of doing so.''
\item Explanation: if a Dalit person is present and I merely think of a casteist remark without uttering it, I am not guilty under law (no right violated) but I am guilty in ethics --- because I am refraining only out of fear of punishment (e.g. the Prevention of Atrocities Act against SC/ST), not out of respect for the moral law. Had I truly respected the law, the thought would not have arisen (just as, respecting our parents, abuse does not come to mind). Kant stresses intentions strongly.
\end{itemize}
\end{PYQLink}

\subsection{The Moral Law \& the Categorical Imperative}

\begin{KeyConcept}
Since acting under outside dictation would violate free will, the moral law is something each individual must generate on their own. As per Kant, the moral law is given to individuals by themselves --- the decision of right and wrong is made by the individual by thinking in a free, rational and logical way.
\end{KeyConcept}

\paragraph{(a) It must become a Universal Law (Universalisability)}
\begin{itemize}
\item Because every individual thinks freely and logically without prejudice, everyone arrives at the same law --- there is no difference between what I think and what you think. Hence the moral law a person gives themselves should be able to become a universal law.
\item Kant's formula: ``Act in such a way that you could will that the maxim of your act become a universal law.'' (`Maxim' = as if everyone acted as you do.)
\end{itemize}

\begin{CriticalPoint}
\begin{itemize}
\item The test of universality exposes the rationality of an argument --- whether we are being selfish, prejudiced or narrow.
\item Sunil seeking a transfer cannot become a universal law: it is not rational to hold that every officer should seek a transfer whenever threatened (if every officer, doctor or engineer ran away, who would run the country?). So the action is biased, not rational.
\end{itemize}
\end{CriticalPoint}

\paragraph{Further tests of universality}

\begin{ExampleBox}
\begin{itemize}
\item Encounter `justice': In a Gurugram case, a child was found murdered in a school and allegations fell on the bus driver; police held a conference and the media built an atmosphere against him. Later it emerged that another school-going child had committed the act. Had a fake encounter been carried out on the driver, justice would never have been served --- so `encounter' cannot become a universal law.
\item Instant justice: In a crowded metro a girl beside you screams that she was inappropriately touched (perhaps due to a jerk). If bystanders also believe in instant justice, they start beating you without any enquiry. You would then demand a fair trial --- and if you want a fair trial for yourself, you must want it for an accused too. Universality exposes this inconsistency.
\item `Encounter after thorough investigation': an officer who performs 49 encounters in six months and becomes a media celebrity cannot deliver the `50th' once investigation shows the accused is not guilty --- conditionality means the rule cannot become universal, and the negative consequences are very high.
\end{itemize}
\end{ExampleBox}

\begin{CriticalPoint}
`An under-the-table bribe is ethical because it speeds up decision-making, like grease in the system' --- the test of universality shows this is neither logical nor rational, so it cannot be a universal law. An illegal act cannot even be considered for universalisation.
\end{CriticalPoint}

\paragraph{(b) Humanity as an End, Never Merely as a Means}

\begin{KeyConcept}
\begin{itemize}
\item ``Act in such a way that you treat humanity never merely as a means but always as an end in itself.'' Because every person is a rational being, we must never use anyone merely as a means; we must treat them as an end --- respecting their dignity. `Merely' means `only.'
\item `Merely' vs respect: if students paid for a course only to get marks and saw the teacher merely as a means to earn marks, that would be `merely as a means.' But where there is mutual respect --- the teacher respecting the student and the student the teacher as human beings --- it is not `merely.'
\item `Whether in your own person or in the person of another': we must not treat even ourselves merely as a means. An aspirant who, after several failed attempts, feels life has no value beyond clearing the exam and dies by suicide is treating their own person merely as a means --- as if life has no other value.
\end{itemize}
\end{KeyConcept}

\begin{ExampleBox}
\begin{itemize}
\item Cadre-management marriage --- two scenarios. Scenario 1: a man dissatisfied with his cadre talks love to a girl with the sole purpose of marrying to change his cadre --- using her merely as a means (the chief intention is exploitation). Scenario 2: two officers of different services genuinely fall in love and then marry and change cadre --- done out of love, not merely as a means. Marriage should happen out of love.
\item The `white man's burden' to `civilise' colonies treated colonised peoples as not rational enough, and ultimately used them merely as a means. Racism and slavery followed the claim ``we are intellectually superior and should decide for the blacks.''
\item A fake cardiologist `treating' a patient with no intention or skill to cure, only to earn money, uses the patient merely as a means --- unlike a genuine cardiologist who treats and charges a fee.
\end{itemize}
\end{ExampleBox}

\begin{DataFact}
\begin{itemize}
\item In manual scavenging there is roughly 97\% Dalit representation, because Dalits are treated merely as a means and not as ends in themselves.
\item A non-profit organisation in Jaipur district, Rajasthan, allegedly sold over 1,500 girls and women from one river section into marriages over eight years --- using human beings merely as a means.
\end{itemize}
\end{DataFact}

\begin{itemize}
\item Other examples of `merely as a means': misleading advertisements / wrong information, bonded labour (not paying workers fully), human trafficking, the illegal sale of organs, brainwashing / terrorism (arming children who cannot consent), using voters only for votes (bribing voters), marriage for dowry, fake relationships, fake ingredients (fake paneer, fake medicine).
\end{itemize}

\paragraph{(c) Legislator and Legislated (Reciprocity)}

\begin{KeyConcept}
``Always act in such a way that you are both a legislator and legislated.'' When you make a law, you must not merely want everyone else to follow it --- you must follow the same law yourself. This is the requirement of reciprocity, and it is what fairness demands.
\end{KeyConcept}

\begin{ExampleBox}
Reservation: I make the law that reservation should exist because there is inequality in society (as a legislator). But if I am born into a higher caste and then say reservation should not exist because it harms me, I have failed to be both legislator and legislated. Fairness demands I also follow the law I make for others.
\end{ExampleBox}

In such a society, human beings are seen as ends whose dignity is respected --- Kant's idea of a `Kingdom of Ends.'

\subsection{Categorical Imperative, Merits \& Demerits of Kant}

\begin{KeyConcept}
\begin{itemize}
\item The moral law must be self-given, universal, and absolute (unconditional). Kant is deadly against conditional laws.
\item Conditional test --- `Lying is right if you have to save your friend's life': imagine a society where everyone believes this. If Suresh and Ramesh are friends and I am chasing Suresh to kill him, Ramesh's lie about where Suresh went will not work on me --- I already know he is lying to save his friend. For a lie to work, the precondition is that I assume he speaks the truth; but here the universal law is lying, so the law is irrational and can never work.
\item Hence such a law must be a categorical imperative --- universal and absolute --- not a hypothetical (conditional) imperative. `Imperative' means `we ought to do this.'
\end{itemize}
\end{KeyConcept}

\paragraph{Merits of Kant}
\begin{itemize}
\item Respect for human dignity.
\item Emphasis on universality and rationality (promotes fact-based, logical decision-making).
\item Avoids the slippery slope. A slippery slope begins when a universal rule is first broken conditionally --- e.g. ``I will study 10 hours daily,'' then take a break `just today,' then keep slipping; or a PDS officer who breaks the rule for a poor woman without documents and then keeps breaking it (eventually for himself). Because Kant's law has no conditionality, there is no slippery slope.
\item Emphasis on objectivity, rationality and integrity.
\end{itemize}

\paragraph{Demerits of Kant}
\begin{itemize}
\item People are judged not only on principles but also on consequences. When the President of India's car was being taken to hospital in Kanpur and a traffic policeman stopped it citing rules, saying ``I was only following the categorical imperative'' is not enough --- one is also responsible for the consequences.
\item Rigidity can lead to inflexible moral rules.
\item There may be a conflict of duties (e.g. RTI Act versus the Official Secrets Act).
\item Kant ignores the diversity of rational opinion --- he assumed everyone reaches the same conclusion, but rational people can genuinely differ.
\end{itemize}

\begin{QuickRevision}
\begin{itemize}
\item Goodwill = the free and rational choice to act according to the moral law, out of respect for that law, not for personal desire, emotion, reward or consequence.
\item Categorical imperative: (i) act on a maxim that can become a universal law; (ii) treat humanity never merely as a means but as an end; (iii) act as both legislator and legislated. The moral law is self-given, universal and absolute.
\item Merits: dignity, universality / rationality, avoids slippery slope, objectivity \& integrity. Demerits: ignores consequences, rigidity, conflict of duties, ignores diversity of rational opinion.
\end{itemize}
\end{QuickRevision}

\section{Teleological Ethics: Utilitarianism}

\subsection{Utilitarianism --- The Core Idea}

\begin{KeyConcept}
\begin{itemize}
\item In deontology an action is ethical if we do our duty / the means are right; in teleology an action is ethical if the consequences / ends are right.
\item Utilitarianism is a school of thought under teleology. `Utility' means usefulness --- the benefit derived as compared to the cost.
\end{itemize}
\end{KeyConcept}

\subsection{Jeremy Bentham --- Quantitative Utilitarianism}

\begin{KeyConcept}
\begin{itemize}
\item ``Nature has placed mankind under the governance of two sovereign masters (pleasure and pain).'' Every human being experiences pleasure and pain.
\item ``An action is moral, or has utility, if it maximises pleasure and minimises pain.''
\end{itemize}
\end{KeyConcept}

\paragraph{Hedonic Calculus}
\begin{itemize}
\item To measure pleasure as well as pain, Bentham gave the Hedonic Calculus. It consists of various parameters --- duration, intensity, certainty, purity, etc. Other things being equal, a longer duration, higher intensity, greater certainty and greater purity mean more pleasure, and so a more ethical action.
\end{itemize}

\begin{CriticalPoint}
Based on this calculus, pleasure and pain can be quantified. According to Bentham there is NO qualitative difference between different kinds of pleasure --- e.g. intellectual pleasures are the same as spiritual pleasures; only quantity matters (``all intoxication is the same intoxication'').
\end{CriticalPoint}

\paragraph{The four sanctions} According to Bentham our actions are regulated by four sanctions (a sanction being a source of pain).

\begin{center}
\begin{tabularx}{\textwidth}{>{\raggedright\arraybackslash}p{2.6cm} >{\raggedright\arraybackslash}p{3.4cm} Y}
\rowcolor{AccentDark}
\thead{Sanction} & \thead{Underlying fear} & \thead{Note} \\
\toprule
Natural & Fear of death & Self-preservation \\
Political & Fear of punishment & Fear of the law \\
Social & Fear of boycott & e.g. the incest taboo --- relationships society does not accept \\
Religious & Fear of God & --- \\
\bottomrule
\end{tabularx}
\end{center}

\begin{KeyConcept}
It is because of these sanctions (pain) that an individual looks beyond his personal pleasures and acts for the happiness of all.
\end{KeyConcept}

\subsection{John Stuart Mill --- Qualitative Utilitarianism}

\begin{KeyConcept}
\begin{itemize}
\item Certain pleasures are qualitatively superior to others --- intellectual and spiritual pleasures are superior to bodily / sensual pleasures.
\item A human being has an inherent sense of dignity with which to test the quality of a pleasure.
\end{itemize}
\end{KeyConcept}

\begin{CriticalPoint}
Key line: ``It is better to be a human being dissatisfied than a pig satisfied.'' A person would rather remain dissatisfied (e.g. not qualify by unfair means) than behave like a pig and get satisfied. A rapist, or an alcoholic lost in bodily pleasure while neglecting family and duty, behaves like a pig --- lacking the inherent sense of dignity that would reveal how the act violates human dignity.
\end{CriticalPoint}

\begin{CriticalPoint}
Mill added a fifth sanction to Bentham's four --- CONSCIENCE. A person avoids wrong actions because their conscience knows what they did (guilt), regardless of belief in God or society. Conscience thus becomes a source of pain that should guide us. (Related ideas: the marshmallow experiment and delayed gratification --- choosing a higher-order pleasure such as clearing the exam over a lower-order pleasure such as a party.)
\end{CriticalPoint}

\subsection{Act Utilitarianism vs Rule Utilitarianism}

\begin{center}
\begin{tabularx}{\textwidth}{>{\raggedright\arraybackslash}p{2.4cm} Y Y}
\rowcolor{AccentDark}
\thead{Aspect} & \thead{Act Utilitarianism (Bentham)} & \thead{Rule Utilitarianism (Mill)} \\
\toprule
What is judged & The rightness of an action is judged by the consequences of the action itself. & We follow rules that generally maximise happiness. \\
Method & Every action is evaluated individually, each time. & Rather than evaluate each action, follow a good general rule. \\
Result & Can be unstable --- e.g. a Robin Hood recalculating every act of looting and distributing. & Provides a more practical and stable ethical guide. \\
\bottomrule
\end{tabularx}
\end{center}

\begin{GovtPolicy}
Liberty example: unrestrained (absolute) liberty for all leads to chaos --- `matsya nyaya' (the big fish eats the small fish) and autocracy --- so overall liberty falls. Liberty with reasonable restrictions maximises overall liberty and yields the greatest good to the greatest number. Article 19 (freedom of expression with reasonable restrictions) is chosen precisely because its consequences are better --- a rule-utilitarian choice.
\end{GovtPolicy}

\begin{CriticalPoint}
Rule utilitarianism vs deontology: both follow rules, but for different reasons. In deontology we follow the moral law for the respect of the moral law (not for consequences); in rule utilitarianism we follow the rule because it leads to greater utility --- the greatest good to the greatest number.
\end{CriticalPoint}

\subsection{Utilitarianism in Governance}

\begin{KeyConcept}
Every governance action carries both merits and demerits; the aim is to maximise net utility --- maximise merit / pleasure and minimise demerit / pain --- to serve the greatest good to the greatest number.
\end{KeyConcept}

\begin{GovtPolicy}
\begin{itemize}
\item Lockdown: merit --- saving lives; demerit --- impact on GDP growth, employment and industry.
\item Dam / development projects: merit --- power generation, irrigation, employment, GDP growth; demerit --- ecological damage and displacement. The Environmental Impact Assessment (EIA) prioritises merit over demerit --- an application of utilitarianism.
\item AFSPA: demerit --- curtails liberty; merit --- saving lives and maintaining law and order.
\end{itemize}
\end{GovtPolicy}

\subsection{Merits \& Demerits of Utilitarianism}

\paragraph{Merits}
\begin{itemize}
\item Result-oriented --- ultimately results matter, which keeps ethics relevant.
\item Practical and action-oriented.
\item Flexible --- rules can change with consequences (unlike Kant's rigidity).
\item A moral form of democracy --- like democracy, it seeks the greater good of the greater number (`morality in the form of democracy').
\item Avoids conflicting rules. For RTI vs the Official Secrets Act: the RTI Act provides that where there is a conflict, information may be disclosed even if protected under the OSA, provided it serves the public interest --- i.e. the greatest good to the greatest number.
\end{itemize}

\paragraph{Demerits}
\begin{itemize}
\item May lead to majoritarianism and the violation of the human rights of minorities --- `greatest good' can be used as a euphemism / shield to justify oppression.
\end{itemize}

\begin{ExampleBox}
\begin{itemize}
\item Nazi Germany justified killing Jews `for the rebirth and greatest good of Germany' --- `greatest good' used as a euphemism to violate the rights of a minority.
\item Social injustice: development projects that displace tribals are justified by utilitarianism, but the costs fall disproportionately on the tribals --- violating their rights.
\end{itemize}
\end{ExampleBox}

\begin{itemize}
\item May promote a slippery slope because of its conditionality (rules keep changing, e.g. lockdowns).
\item Difficulty in measuring happiness --- how, and whose, happiness is measured? (A tribal may be happiest remaining in the forest.)
\end{itemize}

\begin{QuickRevision}
\begin{itemize}
\item Utilitarianism (teleology): an action is ethical if it maximises net utility --- the greatest good to the greatest number.
\item Bentham (quantitative): two sovereign masters (pleasure \& pain), hedonic calculus, no qualitative difference, four sanctions. Mill (qualitative): higher vs lower pleasures, dignity, `better a human dissatisfied than a pig satisfied,' fifth sanction = conscience.
\item Act utilitarianism (Bentham) judges each action; rule utilitarianism (Mill) follows happiness-maximising rules. Merits: relevant, practical, flexible, democratic, avoids conflict. Demerits: majoritarianism, social injustice, slippery slope, hard to measure happiness.
\end{itemize}
\end{QuickRevision}
